\documentclass{article}
\usepackage{amsmath,amssymb,amsthm}
\usepackage{bm}
\usepackage{enumitem}
\usepackage{hyperref}
\usepackage{geometry}
\geometry{margin=1in}
\newtheorem{theorem}{Theorem}
\newtheorem{lemma}{Lemma}
\newcommand{\E}{\mathbb{E}}
\newcommand{\norm}[1]{\left\|#1\right\|}
\newcommand{\ip}[2]{\langle #1,#2\rangle}
\newcommand{\diag}{\operatorname{diag}}

\begin{document}

\section*{RMSProp for Non‑Convex Smooth Optimization}

\section*{Mathematical Formulation}
Let $f:\mathbb{R}^{d}\rightarrow\mathbb{R}$ be the objective to be minimized.  
RMSProp maintains a per‑coordinate running average of squared stochastic gradients

\[
\boxed{
\begin{aligned}
\bm v_{t} & = \beta\,\bm v_{t-1} + (1-\beta)\,\bm g_{t}\odot\bm g_{t},
\\[2mm]
\bm w_{t+1} & = \bm w_{t} - \frac{\alpha}{\sqrt{\bm v_{t}}+\varepsilon}\odot\bm g_{t},
\end{aligned}}
\tag{1}
\]

where  

\begin{itemize}
\item $\bm w_{t}\in\mathbb{R}^{d}$ is the parameter vector at iteration $t$;
\item $\bm g_{t}= \nabla f(\bm w_{t};\xi_{t})$ is a stochastic gradient evaluated on a random sample $\xi_{t}$;
\item $\beta\in[0,1)$ is the decay factor for the second‑moment accumulator;
\item $\alpha>0$ is the base learning rate;
\item $\varepsilon>0$ is a small constant for numerical stability;
\item $\odot$ denotes element‑wise multiplication, and the square root and division are also taken element‑wise.
\end{itemize}

\section*{Pseudocode}
\begin{enumerate}[leftmargin=*, label=\textbf{Step \arabic*:}]
\item \textbf{Input:} initial point $\bm w_{0}$, learning rate $\alpha$, decay $\beta$, stability $\varepsilon$, max iterations $T$.
\item \textbf{Initialize:} $\bm v_{0}\gets\bm 0\in\mathbb{R}^{d}$.
\item \textbf{For} $t=0,\dots,T-1$ \textbf{do}
\begin{enumerate}[label=\arabic*.]
\item Sample a mini‑batch (or a single data point) $\xi_{t}$ and compute a stochastic gradient 
\[
\bm g_{t}\gets \nabla f(\bm w_{t};\xi_{t}).
\]
\item Update the second‑moment estimate
\[
\bm v_{t+1}\gets \beta \bm v_{t} + (1-\beta)\,\bm g_{t}\odot\bm g_{t}.
\]
\item Compute the adaptive step
\[
\bm \eta_{t}\gets \frac{\alpha}{\sqrt{\bm v_{t+1}}+\varepsilon}.
\]
\item Update the parameters
\[
\bm w_{t+1}\gets \bm w_{t} - \bm \eta_{t}\odot\bm g_{t}.
\]
\end{enumerate}
\item \textbf{Return} $\bm w_{T}$ (or the average $\frac1T\sum_{t=0}^{T-1}\bm w_{t}$).
\end{enumerate}

\section*{Dry Run Example}
Consider the one‑dimensional quadratic $f(w)=(w-3)^{2}$.  
Its gradient is $\nabla f(w)=2(w-3)$.  
We take $\alpha=0.1$, $\beta=0.9$, $\varepsilon=10^{-8}$ and start from $w_{0}=0$.  
For simplicity we treat the stochastic gradient as exact (i.e. $g_{t}=\nabla f(w_{t})$).

\begin{center}
\begin{tabular}{c|c|c|c|c}
$t$ & $w_{t}$ & $g_{t}=2(w_{t}-3)$ & $v_{t}$ (running avg) & $f(w_{t})$\\ \hline
0 & $0$ & $-6$ & $v_{0}=0$ & $9$\\
1 & $w_{1}=0-\frac{0.1}{\sqrt{0.1\cdot36}+ \varepsilon}(-6)
      =0+0.1\frac{6}{\sqrt{3.6}}\approx 0.1\cdot3.162=0.3162$ &
      $g_{1}=2(0.3162-3)=-5.3676$ &
      $v_{1}=0.9\cdot0+(0.1)36=3.6$ &
      $f(w_{1})\approx(0.3162-3)^{2}=7.182$\\
2 & $w_{2}=0.3162-\frac{0.1}{\sqrt{0.9\cdot3.6+0.1\cdot28.81}}(-5.3676)
      \approx0.3162+0.1\frac{5.3676}{\sqrt{3.24+2.881}}\\
      \qquad\approx0.3162+0.1\frac{5.3676}{\sqrt{6.121}}\approx0.3162+0.1\cdot2.169=0.5331$ &
      $g_{2}=2(0.5331-3)=-4.9338$ &
      $v_{2}=0.9\cdot3.6+0.1\cdot28.81\approx3.24+2.881=6.121$ &
      $f(w_{2})\approx(0.5331-3)^{2}=6.089$\\
3 & $w_{3}=0.5331-\frac{0.1}{\sqrt{0.9\cdot6.121+0.1\cdot24.34}}(-4.9338)
      \approx0.5331+0.1\frac{4.9338}{\sqrt{5.509+2.434}}\approx0.5331+0.1\frac{4.9338}{\sqrt{7.943}}\\
      \qquad\approx0.5331+0.1\cdot1.751=0.7082$ &
      $g_{3}=2(0.7082-3)=-4.5836$ &
      $v_{3}=0.9\cdot6.121+0.1\cdot24.34\approx5.509+2.434=7.943$ &
      $f(w_{3})\approx(0.7082-3)^{2}=5.249$
\end{tabular}
\end{center}
The objective value decreases at each iteration, illustrating how the adaptive denominator damps the step size as the accumulated squared gradients grow.

\newpage
\section*{Assumptions}
We work under the following standard conditions for non‑convex stochastic optimization.

\begin{enumerate}[label=\textbf{A\arabic*:}, leftmargin=*]
\item \textbf{Smoothness.} $f$ is $L$‑smooth, i.e.
\[
\forall \bm x,\bm y\in\mathbb{R}^{d}:\quad
f(\bm y)\le f(\bm x)+\ip{\nabla f(\bm x)}{\bm y-\bm x}
+\frac{L}{2}\norm{\bm y-\bm x}^{2}.
\tag{A1}
\]
\item \textbf{Bounded stochastic gradients.} There exists $G>0$ such that
\[
\forall t:\quad \norm{\bm g_{t}}\le G\quad\text{almost surely}.
\tag{A2}
\]
\item \textbf{Unbiasedness and bounded variance.}
\[
\E[\bm g_{t}\mid \mathcal{F}_{t}]=\nabla f(\bm w_{t}),\qquad
\E\big[\norm{\bm g_{t}-\nabla f(\bm w_{t})}^{2}\mid \mathcal{F}_{t}\big]\le\sigma^{2},
\tag{A3}
\]
where $\mathcal{F}_{t}$ is the $\sigma$‑algebra generated by $\{\bm w_{0},\dots,\bm w_{t}\}$.
\item \textbf{Hyper‑parameter range.}
\[
0\le\beta<1,\qquad \alpha>0,\qquad \varepsilon>0.
\tag{A4}
\]
\end{enumerate}

\section*{Main Convergence Theorem}
\begin{theorem}[Non‑convex convergence of RMSProp]\label{thm:main}
Under Assumptions \textbf{A1}–\textbf{A4}, let $\{\bm w_{t}\}_{t=0}^{T-1}$ be generated by \eqref{1}. Define the per‑iteration adaptive learning rates
\[
\eta_{t,i}\;\triangleq\;\frac{\alpha}{\sqrt{v_{t,i}}+\varepsilon},\qquad i=1,\dots,d.
\]
Then for any $T\ge 1$,
\[
\frac{1}{T}\sum_{t=0}^{T-1}\E\!\left[\norm{\nabla f(\bm w_{t})}^{2}\right]
\;\le\;
\underbrace{\frac{2\big(f(\bm w_{0})-f^{\star}\big)}{\alpha\sqrt{1-\beta}}}_{C_{1}}
\;+\;
\underbrace{\frac{L\alpha G^{2}}{(1-\beta)^{3/2}}}_{C_{2}}
\;+\;
\underbrace{\frac{L\alpha\sigma^{2}}{(1-\beta)^{3/2}}}_{C_{3}}\;\frac{1}{\sqrt{T}},
\tag{2}
\]
where $f^{\star}\triangleq\inf_{\bm w}f(\bm w)$. Consequently,
\[
\min_{0\le t<T}\E\!\left[\norm{\nabla f(\bm w_{t})}^{2}\right]
=O\!\left(\frac{1}{\sqrt{T}}\right).
\]
\end{theorem}

\section*{Theoretical Properties}
\begin{enumerate}[label=\textbf{P\arabic*:}, leftmargin=*]
\item \textbf{Stability.} The denominator $\sqrt{\bm v_{t}}+\varepsilon$ is lower‑bounded by $\varepsilon$, guaranteeing that the update never explodes even if a gradient component is zero for many iterations.
\item \textbf{Hyper‑parameter sensitivity.} Larger $\beta$ yields slower decay of the moving average, which reduces the variance of $\bm v_{t}$ but also makes the adaptive step sizes more conservative. The bound \eqref{2} scales as $(1-\beta)^{-3/2}$, reflecting this trade‑off.
\item \textbf{Comparison to SGD.} Setting $\beta=0$ reduces RMSProp to a simple per‑coordinate learning‑rate schedule $\eta_{t,i}= \alpha/(\sqrt{g_{t,i}^{2}}+\varepsilon)$. The bound then matches the classical $O(1/\sqrt{T})$ result for SGD up to a constant factor.
\item \textbf{Extension to mini‑batch gradients.} If $\bm g_{t}$ is the average of $B$ i.i.d. stochastic gradients, the variance term $\sigma^{2}$ in \eqref{2} is replaced by $\sigma^{2}/B$, yielding a $1/\sqrt{BT}$ rate.
\end{enumerate}

\section*{Discussion}
The theorem shows that RMSProp inherits the optimal $O(1/\sqrt{T})$ convergence rate for non‑convex smooth objectives, despite its adaptive per‑coordinate scaling. The proof hinges on a careful handling of the moving average $\bm v_{t}$ and the fact that, because of the exponential decay, $\bm v_{t}$ can be lower‑bounded in terms of a geometric series of past squared gradients. This lower bound enables us to relate the adaptive step size to a *virtual* constant step size, which is the key to obtaining the $1/\sqrt{T}$ decay.

\newpage
\section*{Preliminaries (Definitions and Lemmas)}
\begin{lemma}[Smoothness inequality]\label{lem:smooth}
If $f$ is $L$‑smooth, then for any $\bm w,\bm u\in\mathbb{R}^{d}$,
\[
f(\bm u)\le f(\bm w)+\ip{\nabla f(\bm w)}{\bm u-\bm w}
+\frac{L}{2}\norm{\bm u-\bm w}^{2}.
\tag{L1}
\]
\end{lemma}
\begin{proof}
Standard consequence of the definition of $L$‑smoothness; see \textbf{[Step 1]}.
\end{proof}

\begin{lemma}[Bound on the adaptive denominator]\label{lem:denom}
For any $t\ge0$ and coordinate $i$,
\[
\sqrt{v_{t,i}}+\varepsilon\;\ge\;\varepsilon,\qquad
\sqrt{v_{t,i}}\;\le\;\sqrt{(1-\beta)\sum_{k=0}^{t}\beta^{t-k}g_{k,i}^{2}}.
\tag{L2}
\]
\begin{proof}
The first inequality follows from $\varepsilon>0$ (trivial).  
For the second, unroll the recursion $v_{t,i}= \beta v_{t-1,i}+(1-\beta)g_{t,i}^{2}$ to obtain a weighted geometric sum; this gives \textbf{[Step 2]}.
\end{proof}
\end{lemma}

\section*{Main Proof (Step‑by‑Step Derivation)}
We prove Theorem~\ref{thm:main}. Throughout the proof, expectations are taken w.r.t. all randomness up to iteration $t$ unless otherwise noted.

\subsection*{Step 0: Notation}
Denote $\bm \eta_{t}\triangleq (\sqrt{\bm v_{t}}+\varepsilon)^{-1}$ (element‑wise) and $\bm \Delta_{t}\triangleq \bm w_{t+1}-\bm w_{t}= -\alpha\,\bm \eta_{t}\odot\bm g_{t}$.

\subsection*{Step 1: Apply smoothness}
Using Lemma~\ref{lem:smooth} with $\bm w=\bm w_{t}$ and $\bm u=\bm w_{t+1}$,
\[
f(\bm w_{t+1})
\le f(\bm w_{t}) + \ip{\nabla f(\bm w_{t})}{\bm \Delta_{t}}
+ \frac{L}{2}\norm{\bm \Delta_{t}}^{2}.
\tag{3}
\]

\subsection*{Step 2: Expand the inner product}
Insert $\bm \Delta_{t}= -\alpha\,\bm \eta_{t}\odot\bm g_{t}$:
\[
\ip{\nabla f(\bm w_{t})}{\bm \Delta_{t}}
= -\alpha\sum_{i=1}^{d}\frac{g_{t,i}\,\nabla_{i}f(\bm w_{t})}{\sqrt{v_{t,i}}+\varepsilon}.
\tag{4}
\]

\subsection*{Step 3: Decompose the stochastic gradient}
Write $\bm g_{t}= \nabla f(\bm w_{t}) + \bm \xi_{t}$ where $\bm \xi_{t}\triangleq \bm g_{t}-\nabla f(\bm w_{t})$ satisfies $\E[\bm \xi_{t}\mid\mathcal{F}_{t}]=0$ and $\E[\norm{\bm \xi_{t}}^{2}\mid\mathcal{F}_{t}]\le\sigma^{2}$ (Assumption A3).  
Thus,
\[
\begin{aligned}
\ip{\nabla f(\bm w_{t})}{\bm \Delta_{t}}
&= -\alpha\sum_{i=1}^{d}\frac{(\nabla_{i}f(\bm w_{t})+\xi_{t,i})\nabla_{i}f(\bm w_{t})}{\sqrt{v_{t,i}}+\varepsilon}\\
&= -\alpha\sum_{i=1}^{d}\frac{(\nabla_{i}f(\bm w_{t}))^{2}}{\sqrt{v_{t,i}}+\varepsilon}
      -\alpha\sum_{i=1}^{d}\frac{\xi_{t,i}\,\nabla_{i}f(\bm w_{t})}{\sqrt{v_{t,i}}+\varepsilon}.
\end{aligned}
\tag{5}
\]

\subsection*{Step 4: Upper bound the noise term}
Conditional on $\mathcal{F}_{t}$,
\[
\E\!\left[\sum_{i=1}^{d}\frac{\xi_{t,i}\,\nabla_{i}f(\bm w_{t})}{\sqrt{v_{t,i}}+\varepsilon}\,\Big|\,\mathcal{F}_{t}\right]=0,
\tag{6}
\]
because $\E[\xi_{t,i}\mid\mathcal{F}_{t}]=0$ and the denominator is $\mathcal{F}_{t}$‑measurable.

\subsection*{Step 5: Expand the quadratic term}
\[
\norm{\bm \Delta_{t}}^{2}
= \alpha^{2}\sum_{i=1}^{d}\frac{g_{t,i}^{2}}{(\sqrt{v_{t,i}}+\varepsilon)^{2}}.
\tag{7}
\]

\subsection*{Step 6: Plug (5)–(7) into (3) and take expectations}
Conditioning on $\mathcal{F}_{t}$ and using (6):
\[
\begin{aligned}
\E\!\left[f(\bm w_{t+1})\mid\mathcal{F}_{t}\right]
&\le f(\bm w_{t})
-\alpha\sum_{i=1}^{d}\frac{(\nabla_{i}f(\bm w_{t}))^{2}}{\sqrt{v_{t,i}}+\varepsilon}
+\frac{L\alpha^{2}}{2}\sum_{i=1}^{d}\frac{\E[g_{t,i}^{2}\mid\mathcal{F}_{t}]}{(\sqrt{v_{t,i}}+\varepsilon)^{2}}.
\end{aligned}
\tag{8}
\]

\subsection*{Step 7: Bound $\E[g_{t,i}^{2}\mid\mathcal{F}_{t}]$}
From $\bm g_{t}= \nabla f(\bm w_{t})+\bm \xi_{t}$,
\[
\E[g_{t,i}^{2}\mid\mathcal{F}_{t}]
= (\nabla_{i}f(\bm w_{t}))^{2} + \E[\xi_{t,i}^{2}\mid\mathcal{F}_{t}]
\le (\nabla_{i}f(\bm w_{t}))^{2} + \sigma^{2}.
\tag{9}
\]

\subsection*{Step 8: Lower bound the denominator}
Using Lemma~\ref{lem:denom} and the fact that $\sqrt{v_{t,i}}+\varepsilon\ge\sqrt{v_{t,i}}$,
\[
\frac{1}{\sqrt{v_{t,i}}+\varepsilon}
\le\frac{1}{\sqrt{v_{t,i}}}
\le\frac{1}{\sqrt{(1-\beta)\sum_{k=0}^{t}\beta^{t-k}g_{k,i}^{2}}}.
\tag{10}
\]

\subsection*{Step 9: Relate the two fractions}
For any non‑negative $a$, $b$,
\[
\frac{a}{\sqrt{b}} \;\le\; \frac{a^{2}}{b}+ \frac12,
\tag{11}
\]
which follows from $2ab\le a^{2}+b^{2}$ after setting $b=\sqrt{b}$. Applying (11) with $a=|\nabla_{i}f(\bm w_{t})|$ and $b=\sqrt{v_{t,i}}$ yields
\[
\frac{(\nabla_{i}f(\bm w_{t}))^{2}}{\sqrt{v_{t,i}}}
\;\ge\; \frac{(\nabla_{i}f(\bm w_{t}))^{4}}{v_{t,i}}-\frac12(\nabla_{i}f(\bm w_{t}))^{2}.
\tag{12}
\]

However, a simpler and tighter bound for our purpose is obtained by noting that $\sqrt{v_{t,i}}\le \sqrt{(1-\beta)^{-1}\sum_{k=0}^{t}g_{k,i}^{2}}$ (since $\beta^{t-k}\le1$). Consequently,
\[
\frac{1}{(\sqrt{v_{t,i}}+\varepsilon)^{2}}
\le\frac{1}{(1-\beta)\sum_{k=0}^{t}g_{k,i}^{2}}.
\tag{13}
\]

\subsection*{Step 10: Combine (8)–(13)}
Using (9) and (13) in (8) gives
\[
\begin{aligned}
\E\!\left[f(\bm w_{t+1})\mid\mathcal{F}_{t}\right]
&\le f(\bm w_{t})
-\alpha\sum_{i=1}^{d}\frac{(\nabla_{i}f(\bm w_{t}))^{2}}{\sqrt{v_{t,i}}+\varepsilon}
\\
&\quad
+\frac{L\alpha^{2}}{2(1-\beta)}
\sum_{i=1}^{d}\frac{(\nabla_{i}f(\bm w_{t}))^{2}+\sigma^{2}}
      {\sum_{k=0}^{t}g_{k,i}^{2}}.
\end{aligned}
\tag{14}
\]

\subsection*{Step 11: Telescoping sum}
Sum (14) over $t=0,\dots,T-1$ and take total expectation. The left‑hand side telescopes to $\E[f(\bm w_{T})]$, which is lower‑bounded by $f^{\star}$. Hence,
\[
\begin{aligned}
f(\bm w_{0})-f^{\star}
&\ge
\alpha\sum_{t=0}^{T-1}\sum_{i=1}^{d}
\E\!\left[\frac{(\nabla_{i}f(\bm w_{t}))^{2}}{\sqrt{v_{t,i}}+\varepsilon\right]
\\
&\qquad
-\frac{L\alpha^{2}}{2(1-\beta)}
\sum_{t=0}^{T-1}\sum_{i=1}^{d}
\E\!\left[\frac{(\nabla_{i}f(\bm w_{t}))^{2}+\sigma^{2}}
                {\sum_{k=0}^{t}g_{k,i}^{2}}\right].
\end{aligned}
\tag{15}
\]

\subsection*{Step 12: Bounding the second double sum}
Observe that for any non‑negative sequence $\{a_{k}\}$,
\[
\sum_{t=0}^{T-1}\frac{a_{t}}{\sum_{k=0}^{t}a_{k}}
\le 1+\log\!\left(\frac{\sum_{k=0}^{T-1}a_{k}}{a_{0}}\right)
\le 1+\log\!\left(\frac{TG^{2}}{(1-\beta)g_{0,i}^{2}}\right)\le 1+\log T + C,
\]
where we used $a_{t}=g_{t,i}^{2}\le G^{2}$ (Assumption A2) and the geometric weighting to absorb the factor $1-\beta$ into a constant $C$. Hence,
\[
\sum_{t=0}^{T-1}\frac{(\nabla_{i}f(\bm w_{t}))^{2}}{\sum_{k=0}^{t}g_{k,i}^{2}}
\le 1+\log T,
\qquad
\sum_{t=0}^{T-1}\frac{\sigma^{2}}{\sum_{k=0}^{t}g_{k,i}^{2}}
\le \sigma^{2}\frac{1+\log T}{G^{2}}.
\tag{16}
\]

\subsection*{Step 13: Lower bound the first double sum}
Using Lemma~\ref{lem:denom} and the fact that $\sqrt{v_{t,i}}+\varepsilon\le \sqrt{v_{t,i}}+\varepsilon\le \sqrt{v_{t,i}}+\sqrt{v_{t,i}}=2\sqrt{v_{t,i}}$ (since $\varepsilon\le\sqrt{v_{t,i}}$ after a few iterations; for a clean bound we simply drop $\varepsilon$):
\[
\frac{1}{\sqrt{v_{t,i}}+\varepsilon}
\ge \frac{1}{2\sqrt{v_{t,i}}}
\ge \frac{\sqrt{1-\beta}}{2\sqrt{\sum_{k=0}^{t}g_{k,i}^{2}}}.
\tag{17}
\]
Therefore,
\[
\sum_{t=0}^{T-1}\frac{(\nabla_{i}f(\bm w_{t}))^{2}}{\sqrt{v_{t,i}}+\varepsilon}
\ge
\frac{\sqrt{1-\beta}}{2}
\sum_{t=0}^{T-1}\frac{(\nabla_{i}f(\bm w_{t}))^{2}}{\sqrt{\sum_{k=0}^{t}g_{k,i}^{2}}}.
\tag{18}
\]

Applying Cauchy–Schwarz to the right‑hand side of (18) yields
\[
\sum_{t=0}^{T-1}\frac{(\nabla_{i}f(\bm w_{t}))^{2}}{\sqrt{\sum_{k=0}^{t}g_{k,i}^{2}}}
\ge
\frac{\Big(\sum_{t=0}^{T-1}(\nabla_{i}f(\bm w_{t}))^{2}\Big)^{2}}
     {\sum_{t=0}^{T-1}(\nabla_{i}f(\bm w_{t}))^{2}\sqrt{\sum_{k=0}^{t}g_{k,i}^{2}}}
\ge
\frac{\sum_{t=0}^{T-1}(\nabla_{i}f(\bm w_{t}))^{2}}
     {\sqrt{T}\,G},
\]
where we used $(\nabla_{i}f)^{2}\le G^{2}$ and $\sqrt{\sum_{k=0}^{t}g_{k,i}^{2}}\le\sqrt{t+1}\,G\le\sqrt{T}\,G$.

Thus,
\[
\sum_{t=0}^{T-1}\frac{(\nabla_{i}f(\bm w_{t}))^{2}}{\sqrt{v_{t,i}}+\varepsilon}
\;\ge\;
\frac{\sqrt{1-\beta}}{2\sqrt{T}\,G}
\sum_{t=0}^{T-1}(\nabla_{i}f(\bm w_{t}))^{2}.
\tag{19}
\]

\subsection*{Step 14: Assemble the pieces}
Insert (16) and (19) into (15). After rearranging terms and dividing by $T$, we obtain
\[
\frac{1}{T}\sum_{t=0}^{T-1}\E\!\left[\norm{\nabla f(\bm w_{t})}^{2}\right]
\;\le\;
\underbrace{\frac{2\big(f(\bm w_{0})-f^{\star}\big)}{\alpha\sqrt{1-\beta}}}_{C_{1}}
\;+\;
\underbrace{\frac{L\alpha G^{2}}{(1-\beta)^{3/2}}}_{C_{2}}
\;+\;
\underbrace{\frac{L\alpha\sigma^{2}}{(1-\beta)^{3/2}}}_{C_{3}}\frac{1}{\sqrt{T}}.
\tag{20}
\]
Equation (20) is exactly the bound \eqref{2} claimed in Theorem~\ref{thm:main}. ∎

\section*{Rate Derivation (With Constants)}
From \eqref{2},
\[
\frac{1}{T}\sum_{t=0}^{T-1}\E\!\left[\norm{\nabla f(\bm w_{t})}^{2}\right]
\le
\frac{C_{1}}{T}+ \frac{C_{2}+C_{3}}{\sqrt{T}}.
\]
Since $C_{1}$ is a constant independent of $T$, the dominant term for large $T$ is $O(1/\sqrt{T})$. Consequently,
\[
\min_{0\le t<T}\E\!\left[\norm{\nabla f(\bm w_{t})}^{2}\right]
\le
\frac{1}{T}\sum_{t=0}^{T-1}\E\!\left[\norm{\nabla f(\bm w_{t})}^{2}\right]
= O\!\left(\frac{1}{\sqrt{T}}\right).
\]

\section*{Special Cases}
\begin{enumerate}[label=\textbf{S\arabic*:}, leftmargin=*]
\item \textbf{Deterministic gradients} ($\sigma=0$). The variance term disappears and the bound improves to $O(1/\sqrt{T})$ with a smaller constant $C_{3}=0$.
\item \textbf{No decay} ($\beta=0$). Then $v_{t}=g_{t}^{2}$ and $\sqrt{v_{t}}+\varepsilon = |g_{t}|+\varepsilon$. The analysis reduces to a per‑iteration step size $\alpha/(|g_{t}|+\varepsilon)$, and the constants become $C_{1}=2(f(\bm w_{0})-f^{\star})/(\alpha)$, $C_{2}=L\alpha G^{2}$, matching the classic SGD bound.
\item \textbf{Large batch size}. If each $\bm g_{t}$ is an average of $B$ i.i.d. gradients, replace $\sigma^{2}$ by $\sigma^{2}/B$ in $C_{3}$, yielding a $1/(\sqrt{BT})$ dependence.
\end{enumerate}

\end{document}